\documentclass[10pt,a4paper]{amsart}
\usepackage[margin=19mm]{geometry}
\usepackage{amsmath,amssymb,mathtools}
\usepackage[scr=rsfs,cal=euler]{mathalfa}
\usepackage{xfrac}
\usepackage[unicode,hidelinks,bookmarks=false]{hyperref}
\newcommand{\ie}{i.e.\ }
\newcommand{\cf}{cf.\ }
\newcommand{\resp}{respectively\ }
\newcommand{\Subsection}{\subsection}
\providecommand{\nobreakdash}{\nobreak\hbox{-}}
\setlength{\emergencystretch}{2em}
\multlinegap0pt
\usepackage{amssymb}
\newtheorem{theorem}{Theorem}
\title{The Maximum of $\operatorname{per}(I-A)$ in Odd Order}
\author{Yair Lavi}
\date{Source: arXiv:2608.08933v1 (CC BY 4.0)}
\begin{document}
\maketitle

\noindent\emph{Source and licence.} The Maximum of $\operatorname{per}(I-A)$ in Odd Order. Version arXiv:2608.08933v1 (CC BY 4.0), distributed by arXiv under the Creative Commons Attribution 4.0 International licence, http://creativecommons.org/licenses/by/4.0/, as recorded in the arXiv licence metadata for this exact version and shown on the abstract page https://arxiv.org/abs/2608.08933v1. Full licence notice: \texttt{licenses/arxiv-2608.08933-CC-BY-4.0.txt}.\par\smallskip

\noindent\textit{Editorial scope.} Counted unit: the article's theorem theorem (source0.tex lines 481--490). The article's complete original proof of it is delivered with it (source0.tex lines 492--541, one \texttt{proof} environment of 50 lines). No mathematics is written, repaired or re-derived: the statement and the proof are the article's own lines, in the article's own notation, and no retained line is substituted or re-flowed.  No further source-local context is needed: every cross-reference of every retained line resolves inside the two delivered spans.  Nothing was omitted from any retained span, so the package carries no omission record.  No retained line contains a citation, so the wrapper carries no bibliography and no bibliography entry.  No retained line contains a figure environment, an image-inclusion command or drawing code, so no image is delivered and the package declares no asset dependency.  Editorial formatting only, in addition: an AMS \texttt{amsart} preamble that restates the article's own declarations and macros used by the retained lines, namely the environments \texttt{theorem} on separate counters, together with the standard packages those lines need.  The retained environments print in this document as Theorem 1 in order of appearance, whereas the article prints the same items with the numbering of its own full document.  The complete native source of the article as distributed in the arXiv e-print for this exact version is bundled unchanged as \texttt{sources/207203/source0.tex}.
\begin{theorem}[transpose symmetrization]
For every
$A\in\Omega_{2k+1}$ and $X=(A+A^{\mathsf T})/2$,

\[
m_k(A)-Q_3(A)\le m_k(X)-Q_3(X).
\tag{4.10}
\]

\end{theorem}

\begin{proof}
Here the column-sum condition enters the upper-bound argument:
because $A$ is doubly stochastic, $A^{\mathsf T}$ is row-stochastic and
$X$ is symmetric and doubly stochastic. The quantities $m_k$ and $Q_3$ are
transpose invariant. For $m_k$, apply the first formula in (3.7) to both
$A$ and $A^{\mathsf T}$ and use $C(A)=C(A^{\mathsf T})$; for $Q_3$,
transposition exchanges the two orientations in (4.2). Moreover
$H_T(A)=H_T(A^{\mathsf T})$. We may therefore replace $K_T(A)$ in (4.5)
by its transpose average.

For an edge $\{r,s\}$ outside $T$, write

\[
x_{rs}=\frac{a_{rs}+a_{sr}}2,
\qquad
\kappa_{rs}=\frac{a_{rs}-a_{sr}}2.
\]

Then

\[
a_{rs}a_{sr}=x_{rs}^2-\kappa_{rs}^2,
\qquad
0\le x_{rs}^2-\kappa_{rs}^2\le x_{rs}^2.
\tag{4.11}
\]

Expanding the hafnian in (4.5) over perfect matchings and using Lemma 4.1
expresses the charge as a sum of terms

\[
(\Phi_T-D_T)
\prod_{\{r,s\}\in\nu}(x_{rs}^2-\kappa_{rs}^2).
\tag{4.12}
\]

The following sign distinction is essential. If $\Phi_T-D_T\le0$, the term
in (4.12) is nonpositive and hence at most
$\Phi_T\prod_{\{r,s\}\in\nu}x_{rs}^2$. If $\Phi_T-D_T>0$, first discard
$D_T$ and then use (4.11), obtaining the same upper bound. Summing over
triples and outside perfect matchings gives

\[
m_k(A)-Q_3(A)
\le\sum_{|T|=3}\Phi_T H_T(X)
=m_k(X)-Q_3(X),
\]

where the last equality is (4.5) specialized to the symmetric matrix $X$.
\end{proof}

\end{document}
